% p.68 — الباب الموفي بالثلاثين في العوامل ومتعلقاتها
\mspage{68}
% Layout: heading across the page; then two columns (right column first).  The two-line
% prose introduction stands at the head of the right column (above its table); each
% column is a list «رقم + عامل | أمثلة».  Set as `twocol` with `rightbox`/`leftbox` (not
% rightcol/leftcol: reledpar pairs those line by line and a table counts as one line, so
% the tables could not sit beside each other); the right box holds the prose and then its
% table, the left box its table; tables are scaled to the column width.  The particle names stand in
% a narrow sub-column at the right edge of each column, the examples in the wide part,
% often several examples per particle on successive lines (one row per printed line
% below).  Numbers are circled (1)–(5) in the right column and plain in the left;
% all written as (n).  Printed page number ۶۸ at top left.
\bab{الباب الموفي بالثلاثين في العوامل ومتعلقاتها}
\begin{twocol}
\begin{rightbox}\vspace{0pt}% align the tops of the two boxes (minipage[t] aligns first baselines)
اِعْلَمْ أَنَّ الْعَوَامِلَ مِائَةٌ مِنْهَا حُرُوْفٌ تَجُرُّ الْاِسْمَ وَهِيَ تِسْعَةَ عَشَرَ حَرْفًا
\par
\begin{sarf}{|r|r|}
(1) اَلْبَاءُ & مَرَرْتُ بِزَيْدٍ . بِهِ دَاءٌ . كَتَبْتُ بِالْقَلَمِ\\
 & ظَلَمْتُمْ أَنْفُسَكُمْ بِاتِّخَاذِكُمُ الْعِجْلَ\\
 & اِشْتَرَيْتُ الْفَرَسَ بِسَرْجِهِ\\
 & ذَهَبَ اللّٰهُ بِنُوْرِهِمْ ذَهَبْتُ بِزَيْدٍ\\
 & اِشْتَرَيْتُ الْفَرَسَ بِدِرْهَمٍ\\
 & بِاللّٰهِ لَأَفْعَلَنَّ كَذَا\\
 & اِرْحَمْ بِزَيْدٍ . زَيْدٌ بِالْبَلَدِ\\
 & لَا تُلْقُوْا بِأَيْدِيْكُمْ اِلَى التَّهْلُكَةِ\\
(2) مِنْ & سِرْتُ مِنَ \com{الْبَصَرَةِ}{كذا في الأصل بفتح الصاد، ولعل الصواب: الْبَصْرَةِ} اِلَى الْكُوْفَةِ\\
 & أَخَذْتُ مِنَ الدَّرَاهِمِ خَاتَمٌ مِنْ فِضَّةٍ\\
 & فَاجْتَنِبُوا الرِّجْسَ مِنَ الْأَوْثَانِ\\
 & يَغْفِرْ لَكُمْ مِنْ ذُنُوْبِكُمْ\\
(3) اِلٰى & سِرْتُ مِنَ الْبَصَرَةِ اِلَى الْكُوْفَةِ\\
 & وَلَا تَأْكُلُوْا أَمْوَالَهُمْ اِلٰى أَمْوَالِكُمْ\\
 & فَاغْسِلُوْا وُجُوْهَكُمْ وَأَيْدِيَكُمْ اِلَى الْمَرَافِقِ\\
 & صُمْتُ اِلَى اللَّيْلِ\\
(4) فِيْ & اَلْمَالُ فِي الْكِيْسِ . نَظَرْتُ فِي الْكِتَابِ\\
 & وَلَأُصَلِّبَنَّكُمْ فِيْ جُذُوْعِ النَّخْلِ\\
(5) اَللَّامُ & اَلْجُلُّ لِلْفَرَسِ . جِئْتُكَ لِاِكْرَامِكَ\\
\end{sarf}
\end{rightbox}
\begin{leftbox}\vspace{0pt}%
% TODO: first line of the left column: the first word reads like لِلّٰهِ and the second
% example like رَادِفَ لَكُمْ (Qur'anic رَدِفَ لَكُمْ); both doubtful.
\begin{sarf}{|r|r|}
 & \unclear{لِلّٰهِ} لَا يُؤَخِّرُ الْأَجَلُ . \unclear{رَادِفَ لَكُمْ}\\
 & لَزِمَ الشَّرُّ لِلشَّقَاوَةِ\\
(6) رُبَّ & رُبَّ رَجُلٍ كَرِيْمٍ لَقِيْتُهُ\\
 & وَرُبَّهُ رَجُلًا جَوَادٌ\\
(7) وَاوُ رُبَّ & وَ عَالِمٍ يَعْمَلُ بِعِلْمِهِ\\
(8) عَنْ & رَمَيْتُ السَّهْمَ عَنِ الْقَوْسِ\\
(9) عَلٰى & زَيْدٌ عَلَى السَّطْحِ عَلٰى زَيْدٍ دَيْنٌ\\
 & مَرَرْتُ عَلَيْهِ اِنْ كُنْتُمْ عَلٰى سَفَرٍ\\
(10) اَلْكَافُ & زَيْدٌ كَأَسَدٍ لَيْسَ كَمِثْلِهِ شَيْئٌ\\
(11) مُذْ & مَا رَأَيْتُهُ مُذْ يَوْمِ الْجُمُعَةِ\\
 & مَا رَأَيْتُهُ مُذْ يَوْمَيْنِ\\
(12) مُنْذُ & مَا رَأَيْتُهُ مُنْذُ يَوْمَيْنِ\\
(13) حَتّٰى & نِمْتُ الْبَارِحَةَ حَتَّى الصَّبَاحِ\\
 & سِرْتُ الْبَصَرَةَ حَتَّى السُّوْقِ\\
 & قَرَأْتُ وِرْدِيْ حَتَّى الدُّعَاءِ\\
 & أَكَلْتُ السَّمَكَةَ حَتّٰى رَأْسِهَا\\
 & مَاتَ النَّاسُ حَتَّى الْأَنْبِيَاءِ\\
 & جَاءَ الْحُجَّاجُ حَتَّى الْمُشَاةِ\\
(14) وَاوُ الْقَسَمِ & وَاللّٰهِ لَأَشْرَبَنَّ اللَّبَنَ\\
(15) بَاءُ الْقَسَمِ & بِاللّٰهِ لَأَقْتُلَنَّ زَيْدًا\\
(16) تَاءُ الْقَسَمِ & تَاللّٰهِ لَأَنْصُرَنَّ زَيْدًا\\
\end{sarf}
\end{leftbox}
\end{twocol}
